\documentclass[12pt]{article}
\usepackage[utf8]{inputenc}
\usepackage[T1]{fontenc}
\usepackage{hyperref}
\title{Wendat Language Resources for K-12 Curriculum Development}
\author{Compiled by Claude Sonnet 5 (claude-sonnet-5) for Andr\'as Kornai}
\date{}
\begin{document}
\maketitle

\section{Introduction}
Wendat (Huron-Wendat, ISO 639-3 wdt) is an Iroquoian language of Canada. It has been dormant since the 19th century, but a community revitalization effort in Wendake, Qu\'ebec, has revived it. According to a Government of Canada blog post, teaching began in 2010 with night classes, expanded to a school in 2011 and to a daycare in 2015. This dossier is based on a single round of web searches, and the links have not been checked for current availability.

\section{Primary Educational Resources}
The search results did not show any published textbook or primer. The Government of Canada blog post describes the community classes, school and daycare:
\begin{itemize}
\item Wendat kwatatiahtah! (Let's speak Wendat!): \url{https://our-languages.canada.ca/en/blogue-blog/wendat-kwatatiahtah-eng}
\item The same post at another address: \url{https://nos-langues.canada.ca/en/blogue-blog/wendat-kwatatiahtah-eng}
\end{itemize}

\section{Linguistic and Reference Materials}
\begin{itemize}
\item University of Victoria Wendat dictionary project, about page: \url{https://hcmc.uvic.ca/wendat/en/about.html}. According to the search summary, it was created in response to the need for reconstructed content for learners and teachers.
\item ``A Grammar of Wendat'' is a dissertation project funded by Mellon/ACLS. Its sources are 17th- and 18th-century manuscript dictionaries, grammars and ecclesiastical texts: \url{https://www.acls.org/?p=9553}
\end{itemize}

\section{Religious and Community Resources}
\begin{itemize}
\item Wendat language site (Langue wendat), about page: \url{https://languewendat.com/en/?p=691}
\item Wendat Language Revitalization page: \url{https://dgfv6-shne4-32lls.languewendat.com/en/?p=847}. The site says archival dictionaries, grammars and catechisms are being used to reawaken the language.
\item McGill talk on language revitalization by Megan Lukaniec: \url{https://mcling.blogs.mcgill.ca/language-revitalization-talk-megan-lukaniec/}
\end{itemize}

\section{Digital and Audio Resources}
The search summary mentions an open-access trilingual Wendat--French--English online dictionary at wendatlanguage.com. That address did not appear as a URL in the search results, so it is not cited as a link here. No audio or video resources were identified.

\section{Curriculum Development Context}
The language is being taught in a community school and daycare in Wendake. The school-age curriculum materials that exist could not be identified from these searches. Wyandot, a closely related language of Michigan, Kansas and Oklahoma, is also being revitalized, but no resources for it were examined. Standardizing the orthography was an early step in the revival. Next steps are to contact the Huron-Wendat community language programme directly and to review the university dictionary project.
\end{document}